%% Cover letter to the Editor, Journal of Sports Sciences — DRAFT for the authors to edit and sign.
%% Goes to the editor only (not anonymised). The word count and the numbers of tables and figures are
%% checked by audit/kontroll_tall.py against main.tex and supplement.tex; edit them only via the check.
%% «Both authors have approved» may be sent only after the co-author has approved the final version.

\documentclass[11pt,a4paper]{article}

\usepackage[utf8]{inputenc}
\usepackage[T1]{fontenc}
\usepackage{times}
\usepackage[margin=2.2cm]{geometry}
\usepackage[colorlinks=true,urlcolor=blue]{hyperref}
\pagestyle{empty}
\setlength{\parindent}{0pt}
\setlength{\parskip}{0.55em}
\raggedright

\begin{document}

Atle G.\ Guttormsen\\
School of Economics and Business\\
Norwegian University of Life Sciences (NMBU)\\
P.O.\ Box 5003, 1432 Ås, Norway\\
\href{mailto:atle.guttormsen@nmbu.no}{atle.guttormsen@nmbu.no}

\bigskip
The Editor\\
Journal of Sports Sciences

\bigskip
Dear Editor,

\textbf{Submission of an Original Article: ``How much is the wind worth? Sprints and horizontal jumps from age 10 to senior level''}

We submit the manuscript above for consideration as an Original Article in the \emph{Journal of Sports Sciences}.

Wind is the one weather condition athletics regulates, and the $\pm$2.0 m/s legality rule applies to ten-year-olds as to Olympic finalists. Yet almost all empirical estimates of the wind effect come from elite senior athletes. Using every eligible wind-measured sprint and horizontal-jump result in the Norwegian Athletics Federation's registry from age 10 to senior level, the manuscript estimates the wind effect within athlete and season and shows how it scales with age, ability and event. Per metre per second the effect on sprint times is clearly larger for children than for seniors, mainly because slower athletes are exposed to the wind for longer; headwinds cost more than tailwinds give, by about as much as drag physics predicts on the 100 m; and a legal tailwind is worth more than a tenth of a second on the 100 and 200 m in every age band. The same estimator applied to US collegiate and championship-finalist samples places the elite literature at the fast end of the gradient. The results supply age- and event-specific wind corrections for youth lists, records and talent evaluation, which we believe will interest readers working on performance analysis, measurement and youth sport.

The analysis plan was committed to version control before estimation (the study was not externally preregistered), departures from it are listed in the Methods, and the Methods include the required section on sample size estimation and justification. All estimates were reproduced by an independent implementation in a second programming language. The study is not health research under Norwegian law and required no ethics committee approval; Sikt (the Norwegian Agency for Shared Services in Education and Research) confirmed that the use of the data raises no data-protection concerns. The analysis code will be deposited in a public repository on acceptance.

The main text is about 6,000 words (excluding the abstract, references, tables and figure legends), with six tables, five figures and eight supplementary tables. This is above the journal's guideline of approximately 4,000 words. The additional length comes mainly from the external validation against collegiate and championship samples and from reporting the effects by event, sex and age band, which the correction table requires. We would be glad to shorten the text or move material to the supplementary file if you prefer.

The files are prepared for double-anonymous review: the manuscript and the supplementary file contain no author information, and the title page with author details and declarations is uploaded separately. The manuscript has not been published and is not under consideration elsewhere, and both authors have approved the submission. The authors report no competing interests and received no funding for this work. The registry also underlies other work by the corresponding author on youth performance development and participation (one published study is cited in the manuscript); none of that work examines wind.

Yours sincerely,

\medskip
Atle G.\ Guttormsen (corresponding author), on behalf of both authors

\end{document}
